% !TeX root = ../main.tex
\section{Related Work}
\label{sec:related}

Nguyen et al.~\cite{nguyen2026} compared five tree ensembles across the city (AUC 0.983--0.989) and placed Hai Chau in the high and very-high susceptibility classes, and Hoang et al.~\cite{hoang2025} trained XGBoost on 1,000 flood points from media reports (accuracy 0.93). Both used random train--test splits, which overstate accuracy when nearby samples are correlated~\cite{roberts2017}. Doan et al.~\cite{doan2023} simulated the 2022 rainfall with MIKE Urban and found that riverside manholes flooded when high river levels blocked the outfalls, and Nguyen and Quach~\cite{nguyenquach2024} mapped the October 2022 flooding from Sentinel-1 imagery, mostly outside the dense urban core.

Ha et al.~\cite{ha2021} mapped flash-flood susceptibility along a highway in Hoa Binh, and Sadler et al.~\cite{sadler2018} modeled street flood severity from crowdsourced reports. Yuan et al.~\cite{yuan2023} are closest to this work: random forest predicted flooded roads in Harris County, Texas, from topographic, hydrologic, and precipitation features (AUC 0.86 for Hurricane Harvey), but under random cross-validation. Pan et al.~\cite{pan2026} used a Bayesian network to place road flood sensors in New York City; their work observes flooding on segments, whereas this paper predicts it. Graph neural networks such as GraphSAGE~\cite{hamilton2017} can use the connectivity between segments, but they have not been compared with feature-based models for this task.

Le et al. reviewed optimization models for urban route planning~\cite{atr2026} and investigated the spatiotemporal impacts of route optimization on road networks~\cite{treng2026}. More recently, Le et al.~\cite{le2026} optimized urban bus routes in Da Nang with multi-objective deep reinforcement learning, without hazard costs. Linking segment-level flood probabilities to routing connects these two lines of work, yet no prior study combines road-level segment prediction, spatial block validation, probability calibration, and downstream flood-aware routing evaluation.
